% Original successor text.  The D_*.tex files in this directory are audit extracts,
% not successor proofs.  References use the accompanying CE_references.bib.
\section{The punctured torus: what the fixed-group reduction does and does not imply}

We use continuous homomorphisms throughout.  A free profinite group means
the full completion, and a free pro-$p$ group is explicitly so designated.
For a profinite group $H$, write $H^{(p)}$ for its maximal pro-$p$ quotient
and $O^p(H)$ for the kernel of this quotient.  Cohomology of a profinite
group with finite coefficients means continuous cohomology.

\statement{C.1}{Proved}{A nontrivial pro-odd kernel in a projective group cannot be small}
Let $B$ be a closed subgroup of a free profinite group.  Suppose that
\[
 1\longrightarrow K\longrightarrow B
       \xrightarrow{\chi}\widehat{\mathbb Z}\longrightarrow1
\]
is exact.  For every prime $p$, the group $K^{(p)}$ is either trivial
or an infinitely generated free pro-$p$ group.  Consequently, if $K$ is
pro-odd and nontrivial, it has an infinitely generated free pro-$p$
quotient for some odd prime $p$.  In particular such a $K$ is not
topologically finitely generated.  If $K$ is abelian, then $K=1$,
without any pro-odd assumption.

\begin{proof}
A closed subgroup of a free profinite group has cohomological dimension
at most one at every prime and is projective in the category of profinite
groups.  These implications, including the Sylow criterion for
cohomological dimension, are recalled in Section~1 of
\cite{CEEnglerHaranKochloukovaZalesskii2004}.  The maximal pro-$p$
quotient of a projective profinite group is projective in the pro-$p$
category: lift a finite embedding problem with pro-$p$ target through the
original group; its lift factors through the maximal pro-$p$ quotient.
It is therefore free pro-$p$.  This applies to $K$ and to every closed
subgroup of $B$.

Regard $\mathbb Z_p$ as the $p$-coordinate subgroup of
$\widehat{\mathbb Z}$, with zero at all other primes, and put
$S=\chi^{-1}(\mathbb Z_p)$.  We claim that
\[
 O^p(S)=O^p(K),\qquad
 1\longrightarrow K^{(p)}\longrightarrow S^{(p)}
       \longrightarrow\mathbb Z_p\longrightarrow1.             \tag{C.1.1}
\]
The characteristic subgroup $O^p(K)$ is normal in $S$.  Its quotient
in $S$ is an extension of the pro-$p$ group $K^{(p)}$ by
$\mathbb Z_p$, hence is pro-$p$.  Thus $O^p(S)\subseteq O^p(K)$.
Conversely, restriction to $K$ of every finite $p$-group quotient of
$S$ kills $O^p(K)$, proving the reverse inclusion and the exact
sequence.  This argument proves the needed injectivity on the left;
it does not assume that maximal pro-$p$ quotients preserve arbitrary
short exact sequences.

Both nonzero groups in (C.1.1) are free pro-$p$.  If $K^{(p)}$ were
nontrivial and finitely generated, Theorem~2 of
\cite{CEEnglerHaranKochloukovaZalesskii2004}, applied with
$G=S^{(p)}$ and $N=K^{(p)}$, would give
\[
 \operatorname{vcd}_p(G/N)=\operatorname{cd}_p(G)-1=0.
\]
But $G/N=\mathbb Z_p$ has virtual cohomological dimension one: every
open subgroup is again $\mathbb Z_p$.  This contradiction proves the
first assertion.

If $K$ is pro-odd and nontrivial, it has a nontrivial finite quotient
$D$ of odd order.  The odd-order theorem
\cite{CEFeitThompson1963} says that $D$ is solvable.  A nontrivial
finite solvable group has nontrivial abelianization, hence has a
quotient $C_p$ for an odd prime $p$.  The composite $K\to C_p$
shows that $K^{(p)}\ne1$.  A quotient of a finitely generated profinite
group is finitely generated, proving the stated consequence.

Finally, if $K$ is abelian, each free pro-$p$ quotient $K^{(p)}$ is
abelian and therefore has rank at most one.  The first assertion
forces all these quotients to be trivial.  A nontrivial profinite
abelian group has a nontrivial cyclic prime-order quotient, so $K=1$.
\end{proof}

\statement{C.2}{Refuted}{The coarse structural data of LM4 D.78 do not force its kernel to vanish}
There is a two-generated projective profinite group $B$, embedded as a
retract of $\widehat F_2$, with a split epimorphism
$\chi:B\to\widehat{\mathbb Z}$ such that
\[
 B^{\mathrm{ab}}\cong\widehat{\mathbb Z},\qquad
 \dim_{\mathbb F_2}H^1(B_0,\mathbb F_2)=1
       \quad\text{for every open }B_0\leq B,
\]
but $\ker\chi$ is a nontrivial pro-$3$ group.  Thus projectivity,
these mod-two ranks, and the split pro-odd-kernel description alone
do not imply $\ker\chi=1$.  This is not an example of a torus
involution or a refutation of the actual question C.1 of the prompt.

\begin{proof}
Let $L$ be free pro-$3$ on $x,y$, let
$Z_{(3')}=\prod_{p\ne3}\mathbb Z_p$, and let $z$ be its
topological generator with every coordinate equal to one.  The map
$z\mapsto\sigma$, where $\sigma(x)=x$ and $\sigma(y)=y^{-1}$,
defines a continuous action through the quotient
$Z_{(3')}\to C_2$.  Set
\[
 B=L\rtimes Z_{(3')},\qquad a=xz.
\]
The factors $x$ and $z$ commute.  If $e_3\in\widehat{\mathbb Z}$
is the primary idempotent at three, then
$a^{e_3}=x$ and $a^{1-e_3}=z$.  Consequently $a,y$ generate $B$,
and $A=\overline{\langle a\rangle}$ is isomorphic to
$\widehat{\mathbb Z}$.

The homomorphism $L\to\mathbb Z_3$ with $x\mapsto1,y\mapsto0$
is invariant under $Z_{(3')}$.  Together with the quotient to
$Z_{(3')}$ it defines the continuous split epimorphism
\[
 \chi:B\longrightarrow\mathbb Z_3\times Z_{(3')}
                  =\widehat{\mathbb Z},\qquad \chi(a)=1.
\]
Its kernel is $\ker(L\to\mathbb Z_3)$ and contains $y\ne1$.
In the abelianization the relation $zyz^{-1}=y^{-1}$ kills $y$,
because multiplication by two is invertible in a pro-$3$ abelian
group.  There is no relation on $x$ or $z$.  This proves
$B^{\mathrm{ab}}\cong\mathbb Z_3\times Z_{(3')}$.

For an open subgroup $B_0$ the restriction of the quotient to
$Z_{(3')}$ has pro-$3$ kernel $B_0\cap L$ and an open image
$Z_0\leq Z_{(3')}$.  Every homomorphism $B_0\to C_2$ kills
this kernel.  The open procyclic group $Z_0$ has its two-primary
factor isomorphic to $\mathbb Z_2$.  Hence
\[
 \operatorname{Hom}_{\mathrm{cont}}(B_0,C_2)
 =\operatorname{Hom}_{\mathrm{cont}}(Z_0,C_2)\cong C_2.
\]

A Sylow pro-$3$ subgroup of $B$ is $L$.  For $p\ne3$, a Sylow
pro-$p$ subgroup maps isomorphically to the Sylow subgroup
$\mathbb Z_p$ of the quotient: the kernel is pro-$3$, and
surjectivity on Sylow subgroups follows from the profinite Sylow
theorem.  Thus every Sylow subgroup has cohomological dimension
one at its prime.  The criterion cited in C.1 gives
$\operatorname{cd}B\leq1$, so $B$ is projective.  The epimorphism
$\widehat F_2\to B$ defined by the two generators $a,y$ therefore
has a continuous section.  This section realizes $B$ as a closed
retract of $\widehat F_2$.

The script \path{code/check_kernel_model.py} checks the finite
quotients $(C_{3^n})^2\rtimes C_{2^m}$ for $(n,m)=(1,1),(1,2),(2,2)$.
It checks the group law, the epimorphism $\chi$, the two generators,
the derived group, and the invariant-vector witness.  These finite
checks confirm the stated finite models; projectivity and the
assertion for every open subgroup are proved above, independently
of any finite computation.
\end{proof}

\statement{C.3}{Proved}{The exact finite-coefficient witness for the actual torus kernel}
Let $E=F\rtimes\langle v\rangle$ be the full negative torus
extension of LM4 D.78, and let $T=\langle w\rangle$ be its
noncusp section class.  Suppose that the actual fixed subgroup
$B=C_F(w)$ has the split epimorphism
$\chi:B\to\widehat{\mathbb Z}$ of that row, with specified
section $t\mapsto a^t$, and set $A=\overline{\langle a\rangle}$,
$K=\ker\chi$.  For a finite continuous left $E$-set $X$, put
$Q_X=\mathbb F_2[X^T]$, with the induced $B$-action.  Then
\[
 K=1\quad\Longleftrightarrow\quad
 Q_X^B\longrightarrow Q_X^A\text{ is onto for every such }X.
                                                                  \tag{C.3.1}
\]
If $1\ne k\in K$, one can take $X=E/(ATN)$, where $N$ is
an open normal subgroup of $E$ contained in $F$ and chosen so
that $k\notin ATN$.  The identity-coset basis vector is an explicit
witness to failure of surjectivity.

\begin{proof}
The implication from left to right is immediate from $B=A$.
For the converse, $A$ centralizes $T$, and $A\cap K=1$ by
the specified splitting.  Thus $AT$ is a closed subgroup and
$k\notin AT$.  Compactness separates the point $k$ from $AT$
in some finite quotient of $E$, giving the required $N$.
The subgroup $ATN$ is open.  Its identity coset is fixed by
$A$ and by $T$, so its basis vector belongs to $Q_X^A$.
The element $k$ moves that vector.  It consequently does not
belong to $Q_X^B$, and the map in (C.3.1), being the inclusion
of fixed subspaces, is not onto.
\end{proof}

\statement{C.4}{Open}{The actual kernel in prompt C.1}
In the actual torus setting of C.3, prove or refute $K=1$.
The audited structural conclusions of LM4 D.78 imply that $K$ is
pro-odd.  C.1 shows that any counterexample must have an infinitely
generated free pro-$p$ quotient for an odd prime.  It would suffice
to prove that the actual $K$ is finitely generated, or that it is
abelian.  Neither assertion is established here.  C.2 shows exactly
why the abstract projectivity and mod-two rank arguments do not
establish either assertion.  The extra peripheral duality of the
torus extension has not been realized in the countermodel of C.2.

